% Generated by regenerate_supplementary.py from results/features/tables/lv_source_shift.csv
\begin{table}[H]
\centering
\small
\setlength{\tabcolsep}{4pt}
\begin{adjustbox}{max width=\linewidth}
\begin{tabular}{lrr}
\toprule
Feature & Filling & + convex hull \\
\midrule
\texttt{stroke\_volume} & 98.5 & 4.0 \\
\texttt{lv\_ef} & 98.3 & 3.5 \\
\texttt{ratio\_rv\_lv\_es} & 40.4 & 1.8 \\
\texttt{es\_lv\_index} & 28.8 & 1.8 \\
\texttt{es\_lv\_vol} & 28.8 & 1.8 \\
\texttt{es\_thick\_slicestd\_mean} & 16.5 & 7.2 \\
\texttt{es\_thick\_slicemean\_std} & 14.5 & 37.5 \\
\texttt{mass\_to\_edv\_ratio} & 13.8 & 1.9 \\
\texttt{ratio\_rv\_lv\_ed} & 13.8 & 1.9 \\
\texttt{myocardium\_to\_lv\_ed} & 13.8 & 1.9 \\
\bottomrule
\end{tabular}
\end{adjustbox}
\caption{Largest relative change (\%, over the 100 training patients) of a feature when the LV is reconstructed instead of taken from the ground truth. Only 7 patients change at all; the ejection fraction of one of them changes by 98\% with the filling.}
\label{tab:supp-shift}
\end{table}
